\section{Extreme outputs and the exact endpoint}\label{sec:extreme57}
The strict amplitude slack in the polytope construction is not a technical
convenience that can simply be removed. At an extreme output, convex
averaging has no room to hide a different physical state. This observation
gives an exact finite-horizon formula, independently of spectral gap or
synchronization.

\begin{theorem}[The extreme-orbit formula]\label{thm:extreme57}
Let a group act by affine bijections on a compact convex set $Y$. Let
$x_0\in Y$ have an orbit consisting of extreme points of $Y$, and let a
finite nonempty command alphabet act through this group. There is one seed
and one terminal query, with target $u_w x_0$. Put
\[
 \mathcal O_t=\{u_wx_0:|w|=t\},\qquad n_t=|\mathcal O_t|.
\]
Every exact clocked stochastic realization at horizon $N$ satisfies
\begin{equation}\label{eq:extremeprofile57}
                         |S_t|\ge n_t\quad(0\le t\le N).
\end{equation}
One deterministic realization attains all these inequalities simultaneously.
Consequently
\begin{equation}\label{eq:extremewidth57}
                         W_{N,0}=n_N.
\end{equation}
Neither an identity letter nor an exact return is required.
\end{theorem}
\begin{proof}
Fix $t$ and one suffix $v$ of length $N-t$. Such a suffix exists because the
alphabet is nonempty; at $t=N$ use the empty word. Let $D_v(i)\in Y$ be the
conditional expected terminal decoder after this fixed suffix, when the
cut-$t$ label is $i$. Convexity guarantees that it lies in $Y$. For every
prefix $w$ of length $t$, exact correctness says
\[
             u_vu_wx_0=\sum_i p_w(i)D_v(i),
\]
where $p_w$ is the actual distribution at that cut. Extremality forces
$D_v(i)=u_vu_wx_0$ whenever $p_w(i)>0$. Since the action of $u_v$ is
injective, prefixes with different physical outputs have disjoint
nonempty supports in $S_t$. This proves \eqref{eq:extremeprofile57}, including
the terminal cut.

Conversely, at cut $t$ store the current point of $\mathcal O_t$. A command
sends this point to its image in $\mathcal O_{t+1}$; the terminal decoder is
the stored point. The initialization is deterministic. Finally, any one
fixed command gives an injection $\mathcal O_t\to\mathcal O_{t+1}$, so $n_t$
is nondecreasing. The peak of this realization is $n_N$, proving
\eqref{eq:extremewidth57}.
\end{proof}
The same conclusion holds for an orbit on a Euclidean unit sphere with
legal output set its convex hull: every orbit point is extreme by strict
convexity of the ambient Euclidean ball. For density matrices, a rank-one
projection is extreme. Indeed, if $P=\sum_i p_iD_i$ with $D_i\ge0$ and
$\tr D_i=1$, then $\tr((I-P)D_i)=0$ for every positive weight. Positivity
forces $D_i$ to be supported on the range of $P$, and hence $D_i=P$.

\subsection{A fixed nonabelian family with an exponential exact endpoint}
The next example uses two established inputs, explicitly separated from
the realization argument. Breuillard and Gelander prove that a dense
subgroup of a connected semisimple real Lie group contains a dense free
subgroup on two generators \cite[Theorem~1.1]{BreuillardGelander57}; their
later correction supplies repairs and details for other parts of the
original proof \cite{BreuillardGelanderCorrection57}. The spectral-gap input
for dense algebraic unitary generators is \cite{BG2012}. No algorithm for
finding free generators, or for certifying a spectral-gap constant, is
asserted here.

\begin{lemma}[A line with trivial algebraic projective stabilizer]
\label{lem:transcendentalseed57}
Fix a real number $\tau$ transcendental over $\Q$ and put
\[
 v_\tau=\frac{(1,\tau,\ldots,\tau^{d-1})^{\mathsf T}}
                  {(\sum_{j=0}^{d-1}\tau^{2j})^{1/2}},\qquad
 P_\tau=v_\tau v_\tau^*.
\]
An invertible matrix with algebraic entries that preserves the line
$\C v_\tau$ is scalar.
\end{lemma}
\begin{proof}
If $Mv_\tau=\lambda v_\tau$, the eigenvalue $\lambda$ is algebraic because
it is a root of the characteristic polynomial of the algebraic matrix $M$.
Each row of $(M-\lambda I)(1,\tau,\ldots,\tau^{d-1})^{\mathsf T}=0$ is a
polynomial identity in $\tau$ with algebraic coefficients. Transcendence
over $\Q$ implies transcendence over the field of algebraic numbers: an
algebraic relation over a finite number field would make $\tau$ algebraic
over $\Q$ by transitivity of algebraic extensions. Every row polynomial
therefore has all coefficients zero. Thus $M=\lambda I$.
\end{proof}
One may take a fixed computable transcendental such as
$\tau=\sum_{n=1}^\infty10^{-n!}$. For completeness, its truncations have
rational denominator $10^{m!}$ and strictly positive tail at most
$2\,10^{-(m+1)!}$. For a nonzero integer polynomial $A$ of degree $D$, a
rational $p/q$ not among its finitely many roots satisfies
$|A(p/q)|\ge q^{-D}$. If $A(\tau)=0$, a bounded derivative near $\tau$
would instead bound $|A(p/q)|$ by a constant times $q^{-(m+1)}$ along these
truncations, a contradiction for large $m$. This also records precisely
why the chosen seed is not part of the rational-input specialization.

\begin{theorem}[Exponential exact width and polynomial noisy width]
\label{thm:freeendpoint57}
For every $d\ge2$ there is a fixed five-letter algebraic unitary alphabet
\[
                       \{I,U,U^{-1},V,V^{-1}\}\subset SU(d)
\]
and a fixed rank-one seed $P_\tau$ as above such that the full matrix
interface $F(g)=gP_\tau g^*$, with legal decoder set $\mathcal D_d$, has
\begin{equation}\label{eq:freeexact57}
                      W_{N,0}=2\,3^N-1\qquad(N\ge0).
\end{equation}
For each fixed $0<\varepsilon<r_d=\sqrt{1-1/d}$, its Frobenius-error width
instead satisfies
\begin{equation}\label{eq:freenoise57}
 c_{d,\varepsilon}\left(\frac{N}{\log(N+2)}\right)^{d-1}
 \le W_{N,\varepsilon}\le C_{d,\varepsilon}(N+1)^{d-1}.
\end{equation}
At $\varepsilon\ge r_d$ one label suffices. For the same alphabet and seed,
every fixed amplitude $0<a<1$ admits the exact polynomial upper bound in
\eqref{eq:matrixupper57} and the corresponding lower bound in
\eqref{eq:matrixrate57}.
\end{theorem}
\begin{proof}
The subgroup of $SU(d)$ consisting of matrices with algebraic entries is
dense. One direct verification is to approximate a unitary matrix by a
nonsingular matrix with entries in $\Q+i\Q$, apply Gram--Schmidt, and
multiply its final column by the inverse determinant of the resulting
unitary matrix. All entries produced are algebraic; continuity near the
original unitary matrix gives density. Apply
\cite[Theorem~1.1]{BreuillardGelander57} to this subgroup. It contains
algebraic $U,V$ freely generating a dense free group $\Gamma$ of rank two.
The scalar matrices in $SU(d)$ form a finite group. Since a free group is
torsion-free, $\Gamma$ contains no nonidentity scalar matrix.
Lemma~\ref{lem:transcendentalseed57} now shows that the orbit map
$\Gamma\to\mathcal D_d$, $g\mapsto gP_\tau g^*$, is injective.

With the identity available for padding, the products at length $N$ are
exactly the elements of reduced free-word length at most $N$. There are
four choices for the first letter of a nonempty reduced word and three
for each subsequent letter. Thus the ball has cardinality
\[
               1+4\sum_{j=1}^N3^{j-1}=2\,3^N-1.
\]
Apply Theorem~\ref{thm:extreme57} to obtain \eqref{eq:freeexact57}. In fact the
entire pointwise minimum profile is $|S_t|=2\,3^t-1$.

Assign a symmetric positive law to the four nonidentity commands and
positive mass to the identity. The algebraic dense-generator theorem
\cite{BG2012}, with laziness, gives \eqref{eq:groupgap57}. Consequently
Theorem~\ref{thm:matrixwidth57} applies to this same fixed alphabet and
seed. It yields \eqref{eq:freenoise57}, the threshold $r_d$, and the asserted
strict-amplitude conclusion.
\end{proof}
This is a full legal-matrix-output statement. It does not resolve the
zero-error unit-amplitude problem for a single scalar planar readout:
that readout need not expose every orbit point as an extreme decoder
value. Nor does it identify a sampled measurement outcome with a density
matrix stored in a classical label. The distinction between a prescribed
convex output set and a larger simplex is essential here.
